\documentclass{jsarticle}
\usepackage{amsmath}
\usepackage{amssymb}
\begin{document}
\title{３次元多様体のエントロピー量}
\author{Masaaki Yamaguchi}
\date{}
\maketitle

  オイラーの定数は、
  $$ C = \int {1 \over x^s}dx - \log x $$
  $$ = \int \left(\int {1 \over x^s}dx - \log x \right)\mathrm{dvol} $$
  $$ = \int \int{1 \over x^s}\mathrm{dvol} - \int  \log x \mathrm{dvol} $$
  と表されていて、$ \text{ヒッグス場方程式 + オイラーの定数 = ゼータ関数} $
  と求まる。この方程式が、リッチフロー方程式にもなる。このリッチフロー方程式
  は、微分幾何の量子化でもあり、ガンマ関数とベータ関数、
  ガウスの曲面論、ヒッグス場、アインシュタインテンソル、オイラーの定数と
  全部の方程式が入っている。
  この方程式が、相加相乗平均となり、宇宙と原子となり、不確定性原理が
  宇宙の観測系が、宇宙では確率方程式と成り立っているが、異次元が合わさると、
  ノルムとしての３次元多様体の確率分布方程式でもあり、確率ではない,
  非線形方程式となっている。プランク長体積でもあるが、絶対的な方程式
  にもなっている。
  $$ = {d \over df} \int \int{1 \over (y \log y)^{1 \over 2}}dy_m 
  - \int e^{-x}x^{1 - t} \log x dx_m $$
  $$ \int x^{1 - t}e^{-x}dV = \int x^{1 - t}dm, \int x^{1 - t}e^{-x}dV
  = \int x^{1 - t}\mathrm{dvol}, f = \gamma = \Gamma^{'} = \int 
  e^{-x}x^{1 - t} \log x dx $$
  $$ = {d \over d \gamma}\Gamma^{-1} - (\gamma)^{\gamma^{'}} $$
  $$ = e^{-f} - e^{f} $$
  $$ = {d \over df}\int \int{1 \over (y \log y)^{1 \over 2}}dy_m
  - {d \over df}\int\int {1 \over (x \log x)^2}dx_m $$
  $$ = 2 \cos (i x \log x) $$
  $$  {d \over df}\int\int{1 \over (y \log y)^{1 \over 2}}dy_m
  + {d \over df}\int\int{1 \over (x \log x)^2}dx_m $$
  $$ = {d \over df}F = 2 i \sin (i x \log x) $$
  $$ {d \over df}F + \int C dx_m  
  = 2 (\cos( i x \log x) + i \sin(i x \log x)) $$
  $$ = 2 e^{- \theta} = 2 e^{- i x \log x} $$
  $$ = {d \over dt}g_{ij}(t) = - 2 R_{ij} $$
  $$ \log (x \log x) \ge 2 \sqrt{y \log y} $$
  $$ \log(x \log x) = \log x + \log \log x $$
  この宇宙と異次元でのブラックホールとホワイトホールでの、
  宇宙と原子の運動量と位置の絶対的な不確定性原理となり、
  ３次元多様体の確率分布方程式が、運動量と位置エネルギーが
  両方観測できる式にもなっている。
  円周上が固定されているために、運動量と位置エネルギーが
  両方観測できる式になっている。
  カルーツァ・クライン理論とハイゼンベルク方程式が、
  両方入っている。
  $$ \nabla \psi^2 = 8 \pi G\left({p \over c^3} + {V \over S}\right) $$
  $$ \nabla \psi^2 = 8 \pi G \hbar + 8 \pi {V \over S} $$
  ホワイトホールの原子とブラックホールの宇宙のエネルギーのエントロピー量で
  もある。宇宙がブラックホールとしてのヒッグス場方程式として、
  ホワイトホールが原子としてのプランク長体積でもある。
  $$ \square_v = 2 \sqrt{2 \pi G \hbar} + 2 \sqrt{2 \pi {V \over S}} $$
  原子としてのブラックホールのエントロピー量でもある。
  $$ \sqrt{2 \pi T} = 2 \sqrt{2 \pi{V \over S}} $$

  この式から宇宙が無重力であることが表されている。
  $$ 2 \cos (i x \log x) + 2 i \sin(i x \log x) = 2 e^{-f} $$
  $$ 2 \sqrt{2 \pi {V \over S}} = {d \over df} 
  \int \int{1 \over (x \log x)}dx_m$$
  これらの方程式から宇宙と異次元がwarped passegeになっている。:
\end{document}
